\section{Proofs for SPD Compression and Completion}
\label{app:spd-proofs}

\begin{proof}[Proof of \cref{prop:reference-necessary}]
Choose a direct-sum decomposition \(E=W\oplus W'\) with \(W'\ne\{0\}\).
For \(t>0\), let \(T_t\) act as the identity on \(W\) and as multiplication
by \(t\) on \(W'\).  If a selected positive form \(P\) were invariant under
every such map, then for nonzero \(n\in W'\),
\[
  P(n,n)=P(T_tn,T_tn)=t^2P(n,n)
  \qquad\forall t>0,
\]
contradicting \(P(n,n)>0\).  Hence no invariant selector exists without a
reference or equivalent off-visible structure.
\end{proof}

\begin{proof}[Proof of \cref{thm:spd-compression}]
Fix \(P\in\SPD^d\), write \(R=A^\top PA\), and let \(H=H^\top\) be a
tangent at \(P\).  Define
\[
  C=P^{1/2}AR^{-1/2},
  \qquad
  X=P^{-1/2}HP^{-1/2}.
\]
Then \(C^\top C=I_m\) and
\[
  R^{-1/2}A^\top HA R^{-1/2}=C^\top XC.
\]
Complete \(C\) to an orthogonal matrix \([C,N]\).  The exact block identity
is
\begin{equation}
  \fro X^2
  =\fro{C^\top XC}^2
  +2\fro{N^\top XC}^2
  +\fro{N^\top XN}^2.
  \label{eq:block-frobenius-identity}
\end{equation}
Therefore
\[
  \norm{D\pi_A(P)[H]}_{R,\AI}
  =\fro{C^\top XC}
  \le\fro X
  =\norm H_{P,\AI}.
\]
Integrating along curves proves global contraction
\eqref{eq:spd-contraction}.

Congruence by \(P_0^{-1/2}\) in the ambient cone and by \(R_0^{-1/2}\) in
the visible cone reduces the completion problem to
\[
  \min\{\dai(I,X):U_0^\top XU_0=\widetilde R\},
  \qquad
  \widetilde R=R_0^{-1/2}RR_0^{-1/2}.
\]
The candidate
\[
  X_\star=I+U_0(\widetilde R-I)U_0^\top
\]
is positive definite, has compression \(\widetilde R\), and satisfies
\[
  \dai(I,X_\star)=\fro{\log\widetilde R}=\dai(I,\widetilde R).
\]
Contraction proves it is a minimizer.

For uniqueness, suppose feasible \(X\) attains equality and write
\(L=\log X\).  The compressed curve
\(Y_t=U_0^\top\exp(tL)U_0\) joins \(I\) to \(\widetilde R\).  Hence
\[
  \dai(I,\widetilde R)
  \le\operatorname{Len}(Y)
  \le\fro L
  =\dai(I,X)
  =\dai(I,\widetilde R).
\]
Both length inequalities are therefore equalities.  The nonnegative
continuous difference between the constant ambient speed and the compressed
speed has zero integral, so it vanishes for every \(t\), in particular at
\(t=0\).  Equality in \eqref{eq:block-frobenius-identity} there forces
\(N^\top LU_0=0\) and \(N^\top LN=0\).  Thus
\(L=U_0L_0U_0^\top\), and the endpoint constraint gives
\(\exp(L_0)=\widetilde R\).  Therefore \(X=X_\star\).  Undoing congruence
gives \eqref{eq:canonical-completion} and
\eqref{eq:spd-shortest-completion}.

For \(R,S\in\SPD^m\), the normalized lifts are
\(\widetilde R\oplus I\) and \(\widetilde S\oplus I\).  Their ambient AIRM
distance is the visible AIRM distance, and their geodesic keeps the identity
block fixed.  Thus the section is isometric and totally geodesic.  The ball
identity follows by combining contraction with the equal-distance lift.

At a general \(P\), let
\(L_P(h)=D\Psi_{P,A}(A^\top PA)[h]\).  The section identity and isometry give
\[
  D\pi_A(P)[L_P(h)]=h,
  \qquad
  \norm{L_P(h)}_{P,\AI}=\norm h_{A^\top PA,\AI}.
\]
If \(V\in\ker D\pi_A(P)\), then \(L_P(h)+sV\) is another lift of \(h\).
Infinitesimal contraction gives
\(\norm h^2\le\norm{L_P(h)+sV}^2\) for all \(s\), with equality at zero.
Differentiation shows \(L_P(h)\perp V\).  Surjectivity and dimension counting
identify the image of \(L_P\) with the full horizontal space.  This proves
the Riemannian-submersion and horizontal-lift statements.

For an anchor \(P\), write \(\mathcal L_P=\Psi_{P,A}(\SPD^m)\).  The
sections obey the coherence identity
\begin{equation}
  \Psi_{\Psi_{P,A}(S),A}(T)=\Psi_{P,A}(T),
  \qquad S,T\in\SPD^m.
  \label{eq:spd-section-coherence-proof}
\end{equation}
Indeed, congruence and visible-coordinate covariance reduce the identity to
\(P=I\) and an orthonormal channel \(U\).  Then
\(Q=\Psi_{I,U}(S)=I+U(S-I)U^\top\), so \(QU=US\), and the completion formula
gives
\[
  \Psi_{Q,U}(T)
  =Q+QU S^{-1}(T-S)S^{-1}U^\top Q
  =I+U(T-I)U^\top
  =\Psi_{I,U}(T).
\]
Consequently, if \(Q\in\mathcal L_P\), then
\(\mathcal L_Q=\mathcal L_P\).  Differentiating this identity at \(Q\), and
using the horizontal-lift identification above, yields
\[
  T_Q\mathcal L_P
  =\range D\Psi_{Q,A}(A^\top QA)
  =\mathcal H_Q
  \qquad\text{for every }Q\in\mathcal L_P.
\]
Thus through every ambient point there is a connected embedded integral
manifold of the full horizontal distribution.  The distribution is therefore
integrable.  Any horizontal curve beginning at \(P\) remains in
\(\mathcal L_P\) by local uniqueness of integral manifolds; hence every
connected horizontal integral manifold through \(P\) is contained in
\(\mathcal L_P\).  Therefore \(\mathcal L_P\) is the maximal horizontal
leaf.  The isometry and total-geodesy argument above also shows that it is
complete.  The maximal horizontal leaves are thus complete and totally
geodesic, proving that \(\pi_A\) is split-Hadamard.

For the minimax statement, if \(B=0\), the bounded fiber is the singleton
\(\{X_0\}\), and the radius and unique minimax center are immediate.  Suppose
\(B>0\).  Work in the normalized splitting and choose
symmetric \(H\in\R^{(d-m)\times(d-m)}\) with \(\fro H=1\).  The curve
\[
  X_s=\widetilde R\oplus\exp(sH)
\]
stays in the visible fiber and satisfies
\(\dai(X_0,X_s)=|s|\).  This proves infinite diameter.  The two points
\(X_{-B},X_B\) lie in the bounded fiber and are distance \(2B\) apart, so
every center has worst-case distance at least \(B\).  The canonical point
\(X_0\) has worst-case distance at most \(B\).  If another point attained
\(B\), equality in the triangle inequality between the endpoints would make
it their midpoint.  AIRM is uniquely geodesic, so the midpoint is uniquely
\(X_0\).

Finally, for any admissible \((P_0,A)\), every \(P\) with
\(R=A^\top PA\) obeys
\[
  \dai(P_0,P)\ge\dai(A^\top P_0A,R),
\]
and the canonical lift realizes equality.  Monotonicity of \(\rho\) proves
both inequalities in \eqref{eq:spd-decision-reduction}.  The same construction
lifts every reduced minimizer.  Under strict monotonicity, replacing a
noncanonical full point by the strictly nearer canonical point lowers the
objective, proving the final correspondence.
\end{proof}

\begin{proof}[Proof of \cref{cor:visible-update}]
For \(\alpha=Ac\),
\[
  A^\top u_P=-A^\top PA c=-Rc,
\]
so all feasible updates have the same quotient class.

For the converse, use a \(P_0\)-orthogonal splitting of the cotangent space
into the visible span and its complement, and choose coordinates in which
\(\alpha\) has visible coordinate \(c_0\ne0\).  A feasible completion has
blocks
\[
  P=\begin{pmatrix}R&B\\B^\top&C\end{pmatrix}.
\]
The invisible component of \(-P\alpha\) is \(-B^\top c_0\).  Given the
coordinate vector \(z\) of any desired invisible drift, take
\[
  B=-\frac{c_0z^\top}{c_0^\top c_0}.
\]
Then \(-B^\top c_0=z\).  Choosing
\[
  C=B^\top R^{-1}B+tI,
  \qquad t>0,
\]
makes the Schur complement positive, hence \(P\succ0\).  This realizes every
invisible drift.  The final claim follows from the block-diagonal form and
uniqueness of the canonical completion.
\end{proof}

\begin{proof}[Proof of \cref{cor:gaussian-kl}]
Normalize to \(P_0=I\) and an orthogonal frame adapted to \(U_0\).  Every
feasible covariance has blocks
\[
  X=\begin{pmatrix}\widetilde R&C\\C^\top&D\end{pmatrix},
  \qquad
  S=D-C^\top\widetilde R^{-1}C\succ0.
\]
Up to constants independent of \((C,S)\), twice the forward Gaussian KL is
\begin{align*}
  \tr X-\log\det X
  &=\tr\widetilde R-\log\det\widetilde R
  +\tr S-\log\det S
  +\tr(C^\top\widetilde R^{-1}C).
\end{align*}
The last term is uniquely minimized at \(C=0\).  The scalar inequality
\(s-\log s\ge1\) applied to the eigenvalues of \(S\) shows that the middle
term is uniquely minimized at \(S=I\).  Hence
\(X=\widetilde R\oplus I\), the normalized canonical completion.  Congruence
returns \(\Psi_{P_0,A}(R)\).
\end{proof}

\begin{proof}[Proof of \cref{cor:prior-data}]
For any \(P\), put \(R=A^\top PA\).  Canonical completion gives
\[
  \dai(P_0,P)\ge\dai(R_0,R),
\]
with equality exactly at \(P=\Psi_{P_0,A}(R)\).  The full problem therefore
reduces exactly to
\[
  \min_{R\in\SPD^m}
  (1-\tau)\dai(R_0,R)^2+\tau\dai(R,\Rhat)^2.
\]
Let \(x=\dai(R_0,R)\), \(y=\dai(R,\Rhat)\), and
\(D=\dai(R_0,\Rhat)\).  The triangle inequality gives \(x+y\ge D\), while
\[
  (1-\tau)x^2+\tau y^2
  =\tau(1-\tau)(x+y)^2+((1-\tau)x-\tau y)^2.
\]
The lower bound \(\tau(1-\tau)D^2\) is attained only when \(R\) lies on the
unique AIRM geodesic from \(R_0\) to \(\Rhat\) with
\(x=\tau D\), \(y=(1-\tau)D\).  Thus
\(R=R_0\#_\tau\Rhat\), and canonical completion proves
\eqref{eq:prior-data-solution} and the distance identities.

For robustness, geodesic convexity of distance in a Hadamard manifold gives
\[
  \dai(R_0\#_\tau R_1,R_0\#_\tau R_2)
  \le\tau\dai(R_1,R_2).
\]
The completion section is isometric, so the same bound holds for the full
solutions.
\end{proof}
